\section{Gobernanza e interpretabilidad}

\subsection{Drift Playbook}

\begin{figure}[H]
\centering
\includegraphics[width=0.8\textwidth]{slide-02-08.jpg}
\caption{La diapositiva 02-08 muestra el drift response workflow, que incluye alert, human review y rollback.}
\end{figure}

El Drift Playbook incluye: Detect (alert rules), Diagnose (cluster drift prompt), Decide (governance board), Document (runbook entry). Cada respuesta debe ir acompañada de snapshot, la curva de KL y el human override log.

\begin{table}[H]
\centering
\begin{tabular}{p{0.22\textwidth} p{0.65\textwidth}}
\toprule
Etapa & artifact principal \\
\midrule
Detect & Reward drift alert + latency spike dashboard \\
Diagnose & Prompt clustering report + human override video \\
Decide & Governance board minutes + rollback checklist \\
Document & Runbook entry + postmortem draft \\
\bottomrule
\end{tabular}
\caption{Entrega de artifacts del Drift Playbook}
\end{table}

\begin{knowledgebox}{Los criterios de las alertas automatizadas}
Se recomienda usar tres líneas de alerta: KL drift, hallucination flag y human override frequency; la governance review solo se activa cuando las tres exceden el umbral al mismo tiempo, para evitar false positives frecuentes.
\end{knowledgebox}

\subsection{Automated Evaluation Matrix}

Automation coverage incluye unit-level tests (safety prompt injection, allowed domains), system-level evaluation (trajectory success rate) y ops-level metrics (human override latency). Estas métricas deben mostrarse por niveles en el dashboard y someterse a un cross-check con benchmark/human eval.

\begin{importantbox}{El significado multinivel de la evaluation matrix}
Separar benchmark, safety y ops en niveles, y luego representar el estado de salud con un color-coded heatmap, permite que compliance, ops y research compartan una misma percepción de los hechos en una sola review.
\end{importantbox}

\subsection{Supervisión humana y reuniones de gobernanza}

Cada versión debe designar un alignment owner, un ops owner y un legal owner. El alignment owner mantiene la evidence matrix; el ops owner monitorea los drift dashboards; el legal owner decide si un high-risk prompt necesita un guardrail adicional.

\begin{warningbox}{No pase por alto la governance meeting}
Si la governance meeting es solo una formalidad, el equipo terminará alineándose con la línea que le resulte más fácil de cumplir cuando los indicadores entren en conflicto (generalmente el benchmark). Es necesario preparar un evidence scorecard antes de la reunión y publicar los action items después.
\end{warningbox}

\subsection{Resumen de la sección}

El ciclo cerrado de gobernanza e interpretabilidad es una accountability practice: un drift playbook de alta calidad, una automated evaluation matrix y una governance meeting con múltiples roles le dan al deployment de aprendizaje por imitación trazabilidad y capacidad de respuesta rápida.
